\section*{Resumen}
\addcontentsline{toc}{section}{Resumen}
En este informe analizo la arquitectura distribuida de Netflix, centrándome específicamente en su backend de microservicios en AWS y en Open Connect, su red propia para la entrega de video. Para evaluarlo, apliqué los conceptos de la Unidad 1 (transparencias, teorema CAP/PACELC, modelos de comunicación y tolerancia a fallos) y contrasté estas decisiones de diseño con la arquitectura de Spotify. Además, complementé la teoría con un experimento práctico usando Docker Compose, donde desplegué el servidor Eureka con tres copias de un servicio. Al provocar un fallo, comprobé cómo el sistema prioriza la disponibilidad: Eureka retuvo el nodo caído durante unos 30 segundos en su registro antes de eliminarlo, mientras que procesar su recuperación tomó solo unos 8 segundos. A partir de esto, concluyo que la resiliencia en estas arquitecturas es una tarea compartida con el cliente. Si bien los microservicios son la clave para sostener el tráfico global de Netflix, su complejidad y su consistencia eventual representan una sobreingeniería innecesaria para sistemas locales que exigen datos exactos al instante.

\newpage
